% ---------------------------------------------------------------
% Style des blocs de code pour pandoc -> LaTeX -> PDF (tectonic)
% Usage : pandoc doc.md -o doc.pdf --pdf-engine=tectonic \
%           -H style-code.tex --highlight-style=tango
% ---------------------------------------------------------------

\usepackage{fvextra}      % Verbatim amelioré : retour à la ligne automatique
\usepackage{xcolor}
\usepackage{tcolorbox}
\tcbuselibrary{skins,breakable}

% --- Symboles Unicode absents de Latin Modern (⚠, ✅, ❌, ℹ) ---
% fontawesome5 fournit des icônes vectorielles, sans dépendre d'une police système.
\usepackage{fontawesome5}
\usepackage{newunicodechar}
\definecolor{warnorange}{HTML}{E8A33D}
\newunicodechar{⚠}{{\textcolor{warnorange}{\faExclamationTriangle}}}
\newunicodechar{✅}{{\textcolor{green!60!black}{\faCheckCircle}}}
\newunicodechar{❌}{{\textcolor{red!70!black}{\faTimesCircle}}}
\newunicodechar{ℹ}{{\textcolor{blue!60!black}{\faInfoCircle}}}

% --- Palette (style GitHub clair) ---
\definecolor{codebg}{HTML}{F6F8FA}      % fond du bloc
\definecolor{codeframe}{HTML}{D0D7DE}   % bordure
\definecolor{inlinebg}{HTML}{EFF1F3}    % fond du code inline

% --- Bloc de code : boîte arrondie, coupable sur plusieurs pages ---
\renewenvironment{Shaded}{%
  \begin{tcolorbox}[
    breakable,
    enhanced,
    colback=codebg,
    colframe=codeframe,
    boxrule=0.6pt,
    arc=3pt,
    left=8pt, right=8pt, top=6pt, bottom=6pt,
    before skip=10pt, after skip=10pt
  ]%
}{%
  \end{tcolorbox}%
}

% --- Contenu du bloc : police plus petite + wrapping des longues lignes ---
\DefineVerbatimEnvironment{Highlighting}{Verbatim}{%
  commandchars=\\\{\},
  fontsize=\small,
  breaklines,
  breakanywhere,
  breaksymbolleft={\footnotesize\textcolor{gray}{$\hookrightarrow$\space}}
}

% --- Code inline `comme ceci` : petit badge gris ---
\let\oldtexttt\texttt
\renewcommand{\texttt}[1]{%
  \colorbox{inlinebg}{\small\oldtexttt{#1}}%
}
